
\chapter*{Introducción}
\addcontentsline{toc}{chapter}{Introducción}

%\colorbox{yellow}{\underline{\bf Introducción super general de los centros NV}}
% % % % % % % % % % % % % % % % % % % % % % % % % % % %

%\colorbox{orange}{\underline{\bf Intro centros NV} }
%%%%%%%%%%%%%%%%%%%%%%%%%%%%%%%%%%%%%%%%%%%

El estudio de sistemas binarios eclipsantes constituye una herramienta fundamental en astrofísica, ya que permite determinar parámetros físicos esenciales de las estrellas. Sin embargo, la resolución del problema inverso asociado al ajuste de modelos teóricos a datos observacionales representa un desafío computacional importante debido a la complejidad y dimensionalidad del espacio de parámetros involucrado.\\
En este trabajo se propone un modelo simplificado de sistemas binarios eclipsantes basado en el potencial de Roche, considerando componentes aproximadamente esféricos, órbitas circulares y rotación rígida, con el objetivo de reproducir adecuadamente las curvas de luz observacionales y estimar de manera confiable los parámetros físicos fundamentales del sistema. Para ello, se implementa un sistema de ajuste automático mediante algoritmos genéticos y se emplea el método de Monte Carlo para estimar las incertidumbres asociadas a los parámetros obtenidos.\\
La metodología desarrollada comprende la generación de curvas de luz sintéticas y su comparación con curvas observacionales obtenidas a partir de datos fotométricos. El grado de concordancia entre ambas se evalúa mediante el análisis de los residuos, permitiendo identificar el conjunto de parámetros que proporciona el mejor ajuste. Posteriormente, los resultados obtenidos se comparan con valores reportados en la literatura para sistemas previamente estudiados, con el fin de validar la fiabilidad del modelo y del código de optimización implementado. Finalmente, una vez validado el procedimiento, el modelo es aplicado al estudio de un sistema binario eclipsante cuyo comportamiento no ha sido caracterizado previamente, demostrando así la utilidad y el potencial del enfoque computacional propuesto para futuras investigaciones en el campo de la astrofísica estelar.


